% 小熊座（综述）
% license CCBYSA3
% type Wiki

本文根据 CC-BY-SA 协议转载翻译自维基百科\href{https://en.wikipedia.org/wiki/Ursa_Minor}{相关文章}。

小熊座（Ursa Minor，拉丁语意为“较小的熊”，与大熊座 Ursa Major 相对），又称小熊星座，是位于遥远北天的一组星座。与大熊座类似，小熊座的尾部也可被视作一只勺子的把手，因此在北美被称为“小北斗”（Little Dipper）：由七颗恒星组成，其中四颗构成勺状的斗部，形态与其对应的大北斗相似。小熊座是公元$^{[2]}$世纪天文学家托勒密列出的$^{[48]}$个星座之一，至今仍是现代$^{[88]}$个星座之一。由于北极星位于其中，小熊座在传统航海导航中，尤其对水手而言，具有重要意义。

北极星（Polaris）是该星座中最亮的恒星，为一颗黄白色超巨星，同时也是夜空中最亮的造父变星，其视星等在$^{[1.97]}$至$^{[2.00]}$之间变化。小熊座$\beta$星（Beta Ursae Minoris），又名 Kochab，是一颗已进入晚期演化阶段的恒星，因膨胀并冷却而成为一颗橙色巨星，视星等为$^{[2.08]}$，仅略暗于北极星。Kochab 与视星等约为$^{[3]}$等的小熊座$\gamma$星（Gamma Ursae Minoris）常被称为“北极星的守护者”或“极星守卫者”$^{[4]}$。在该星座中已有行星被探测到环绕其中四颗恒星运行，包括 Kochab。此外，小熊座还包含一颗孤立的中子星——Calvera，以及迄今发现的第二高温白矮星 H1504+65，其表面温度约为$^{[200{,}000]}$ K。
\subsection{历史与神话}
\begin{figure}[ht]
\centering
\includegraphics[width=6cm]{./figures/9a220c9be61db089.png}
\caption{天龙座环绕着小熊座，如约 1825 年在伦敦出版的一组星座图《乌拉尼娅之镜》（Urania's Mirror）中所描绘的那样。} \label{fig_Ursa_1}
\end{figure}
在巴比伦的星表中，小熊座被称为“天之车”（MUL\,MAR.GÍD.DA.AN.NA，并与女神 Damkina 相关联）。它被列入约公元前 1000 年编纂的《MUL.APIN》星表之中，属于“恩利尔之星”，即北天星区。$^{[6]}$ 据第欧根尼·拉尔修记载并引述卡利马科斯的话，米利都的泰勒斯“测量了腓尼基人航海所依凭的那辆车的星辰”。第欧根尼将其认定为小熊座；由于据说腓尼基人在海上航行时使用该星座，它也被称为 Phoinikē。$^{[7][8]}$ 将北天星座称作“熊”的传统似乎确实源自希腊，尽管荷马只提到过一只“熊”。$^{[9]}$ 因而最初的“熊”应为大熊座；而小熊座作为第二只熊，或“腓尼基之熊”（Ursa Phoenicia，即 Φοινίκη），则是后来才被接纳的。根据斯特拉波（I.1.6, C3）的说法，这是由于泰勒斯的建议，他向希腊人推荐该星座作为航海指引，而此前希腊人一直以大熊座进行航行定位。在古典时代，天极的位置略更接近小熊座的 $\beta$ 星而非 $\alpha$ 星，因此整个星座都被视为指示北方。自中世纪起，使用小熊座 $\alpha$ 星（即“北极星”）作为北极指向变得更加便利，尽管在中世纪北极星仍与天极相距数度。$^{[10][a]}$ 如今，北极星已位于距离北天极 1° 以内，并仍然是当前的极星。其新拉丁名称 \emph{stella polaris} 仅在近代早期才出现。$^{[4]}$  

该星座的另一古名为 Cynosura（希腊语 Κυνοσούρα，“狗之尾”）。这一名称的起源并不明确（若将小熊座视为“狗尾”，则意味着附近应存在一只“狗”的星座，但并无此星座）。$^{[11]}$ 相反，在《星变志》（Catasterismi）的神话传统中，Cynosura 被解释为一位俄瑞阿得斯山林仙女，她曾担任宙斯的乳母，因而受到宙斯的褒奖，被安置于天穹之中。$^{[12]}$ 关于 Cynosura 这一名称还存在多种解释：有一种将其与卡利斯托的神话联系起来，认为其子阿耳卡斯被替换为她的狗，并由宙斯将其置于天上；$^{[11]}$ 也有人认为在更早的观念中，大熊座被视为一头牛，与牧人座构成牧牛组合，而小熊座则是一只狗。$^{[13]}$ 乔治·威廉·考克斯将该名称解释为 Λυκόσουρα 的变体，理解为“狼之尾”，但他本人又将其词源解释为“光的轨迹或行列”（即 λύκος“狼”与 λύκ-“光”之间的关联）。艾伦则指出，可将其与古爱尔兰语中该星座的名称 drag-blod（“火之轨迹”）进行比较。布朗（1899）提出 Cynosura 可能并非希腊语起源，而是借自亚述语 An-nas-sur-ra，意为“高升者”。$^{[14]}$  

另有一种神话叙述两只熊曾在伊达山上将宙斯隐藏起来，使其免遭其父克洛诺斯的杀害；后来宙斯将它们置于天空之中，但由于被神抛向天穹，它们的尾巴被拉得很长。$^{[15]}$  

由于小熊座由七颗恒星组成，拉丁语中表示“北方”的词 \emph{septentrio}（即北极星所指的方向）源自 \emph{septem}（七）与 \emph{triones}（牛），意为“拉犁的七头牛”，而这七颗星也常被想象为这种形象。该名称同样被用来指称大熊座的主星。$^{[16]}$  

在因纽特天文学中，最亮的三颗星——北极星、Kochab 与 Pherkad——被称为 Nuutuittut（“永不移动者”），尽管这一称呼更常以单数形式使用，专指北极星。由于在极北纬地区北极星在天空中的高度过高，它并不适合用于航海定位。$^{[17]}$ 在中国天文学中，小熊座的主要恒星被划分为两个星官：勾陈（Gòuchén，包含 $\alpha$、$\delta$、$\varepsilon$、$\zeta$、$\eta$、$\theta$、$\lambda$ UMi）以及北极（Běijí，包含 $\beta$ 与 $\gamma$ UMi）。$^{[18]}$  
\subsection{特征}
\begin{figure}[ht]
\centering
\includegraphics[width=8cm]{./figures/a72e9261fedb4cea.png}
\caption{约公元 $^{[1009-1010]}$ 年《定星书》中所描绘的小熊座。与西方的表现方式不同，图中的熊被画成尾巴下垂的形态。} \label{fig_Ursa_2}
\end{figure}
小熊座西侧与鹿豹座相邻，西侧亦与天龙座接壤，东侧则与仙王座相邻。其覆盖面积为 $256$ 平方度，在 $88$ 个星座中按面积大小排名第 $56$ 位。小熊座在美国的通俗称呼为“小北斗”，因为其最亮的七颗恒星似乎构成了一个长柄勺（舀勺）的形状。勺柄末端的恒星是北极星 Polaris。北极星也可以通过沿着大熊座“斗碗”末端形成的两颗恒星——大熊座 $\alpha$ 星与 $\beta$ 星（通常称为“指极星”）连成的直线，在夜空中延伸约 $30^\circ$（相当于手臂伸直时三个竖起的拳头的宽度）来找到 $^{[19]}$。构成小北斗“斗碗”的四颗恒星，其视星等分别为二等、三等、四等和五等星，为判断肉眼可见恒星的星等提供了一个简便的参照，对城市居民或用于测试视力尤为有用 $^{[20]}$。

该星座的三字母缩写由国际天文学联合会（IAU）于 $^{[1922]}$ 年正式采用，为“UMi”$^{[21]}$。官方星座边界由比利时天文学家欧仁·德尔波特于 $^{[1930]}$ 年确定，并以一个由 $22$ 条边组成的多边形来界定（见信息框示意）。在赤道坐标系中，这些边界的赤经范围介于 $08^\mathrm{h}\,41.4^\mathrm{m}$ 与 $22^\mathrm{h}\,54.0^\mathrm{m}$ 之间，而赤纬范围则从北天极延伸至南侧的 $65.40^\circ$ $^{[1]}$。由于其位于遥远的北天半球，小熊座整体仅能被北半球的观测者完整看到 $^{[22][b]}$。
\subsection{特征}
\begin{figure}[ht]
\centering
\includegraphics[width=6cm]{./figures/9675489a80e4d9ed.png}
\caption{肉眼所见的小熊座星座（已添加连线与标注）。注意大熊座中构成北斗七星的七颗恒星，并从北斗最外侧的两颗恒星（有时称为“指极星”）连成一条直线指向北极星 Polaris。} \label{fig_Ursa_3}
\end{figure}
\subsubsection{恒星}  
德国制图学家约翰·拜耳使用希腊字母 $\alpha$ 至 $\theta$ 来标注该星座中最显著的恒星，而他的同乡约翰·埃勒特·博德随后又补充使用了 $\iota$ 至 $\phi$。目前仅有 $\lambda$ 与 $\pi$ 仍在使用，这很可能是由于它们靠近北天极的缘故 $^{[16]}$。在该星座的边界之内，共有 $39$ 颗视星等亮于或等于 $6.5$ 等的恒星 $^{[22][c]}$。  

按照约翰·拜耳的排序，主要七颗恒星的传统名称为：  
\begin{enumerate}
\item Polaris  
\item Kochab  
\item Pherkad  
\item Yildun  
\item Epsilon Ursae Minoris 无传统名称。  
\item Zeta Ursae Minoris 无传统名称。  
\item Eta Ursae Minoris 无传统名称。  
\end{enumerate}
作为小熊尾部的标志 $^{[16]}$，北极星 Polaris（即小熊座 $\alpha$ 星）是该星座中最亮的恒星，其视星等在 $1.97$ 与 $2.00$ 之间变化，周期为 $3.97$ 天 $^{[24]}$。它距离地球约 $432$ 光年 $^{[25]}$，是一颗黄白色超巨星，光谱型在 F7Ib 与 F8Ib 之间变化 $^{[24]}$，其质量约为太阳的 $6$ 倍，光度约为太阳的 $2500$倍，半径约为太阳的$45$ 倍。北极星是从地球可见的最亮的造父变星。它是一个三星系统，其中的超巨星主星拥有两颗黄白色主序星伴星，轨道半长轴分别约为 $17$ 和 $2400$ 天文单位（AU），其公转周期分别为 $29.6$ 年与 $42000$ 年 $^{[26]}$。  

传统名称为 Kochab 的小熊座$\beta$星，其视星等为$2.08$，略暗于北极星$^{[27]}$。它距离地球约$131$光年$^{[28][d]}$，是一颗橙色巨星——即核心氢已耗尽、离开主序阶段的演化恒星——光谱型为 K4III$^{[27]}$。Kochab 在约 $4.6$ 天的周期内呈现轻微的光变，其质量通过这些振荡测量估计约为太阳的$1.3$倍$^{[29]}$。Kochab 的光度约为太阳的 $450$倍，直径约为太阳的$42$倍，表面温度约为$4130\,\mathrm{K}$,$^{[30]}$。其年龄估计约为$29.5$亿年，误差约 $\pm 10$ 亿年，并已被宣布拥有一颗行星伴星，其质量约为木星的$6.1$倍，轨道周期为$522$天$^{[31]}$。
\begin{figure}[ht]
\centering
\includegraphics[width=6cm]{./figures/58bc21206ce95031.png}
\caption{小熊座与大熊座相对于北极星 Polaris 的位置关系} \label{fig_Ursa_4}
\end{figure}
传统名称为 Pherkad 的小熊座$\gamma$星，其视星等在 3.04 与 3.09 之间变化，周期约为 3.4 小时$^{[32]}$。它与 Kochab 一同被称为“北极星的守护者”$^{[33]}$。该星是一颗光谱型为 A3II--III 的白色亮巨星 $^{[32]}$，质量约为太阳的 4.8 倍，光度约为太阳的 1050 倍，半径约为太阳的 $15$ 倍 $^{[34]}$，距离地球$487\pm8$光年$^{[28]}$。Pherkad 属于盾牌座 $\delta$ 型变星这一恒星类别$^{[32]}$——这是一类短周期（最长不超过 6 小时）的脉动变星，被用作标准烛光并用于研究星震学 $^{[35]}$。可能同属该类别的还有小熊座 $\zeta$ 星 $^{[36]}$，这是一颗光谱型为 A3V 的白色恒星$^{[37]}$，已开始冷却、膨胀并增亮。它很可能曾是一颗 B3 型主序星，目前呈现轻微变光$^{[36]}$。在视星等 4.95 时，小熊座 $\eta$ 星是小北斗七星中最暗的一颗$^{[38]}$。它是一颗光谱型为 F5V 的黄白色主序星，距离地球 97 光年$^{[39]}$，直径为太阳的 $2$ 倍，质量为太阳的 1.4 倍，光度为太阳的 7.4 倍$^{[38]}$。其附近的小熊座$\zeta$ 星旁还有视星等为 5.00 的小熊座$\theta$星，该星距离地球$860\pm80$ 光年 $^{[40]}$，是一颗光谱型为 K5III 的橙色巨星，已离开主序阶段并发生膨胀与冷却，其直径估计约为太阳的 4.8 倍 $^{[41]}$。

构成小北斗“斗柄”的恒星为小熊座 $\delta$ 星（Yildun）$^{[42]}$与小熊座 $\varepsilon$ 星。$\delta$ 星距离北天极略超过$3.5^\circ$，是一颗光谱型为 A1V 的白色主序星，视星等为 4.35$^{[43]}$，距离地球$172\pm1$光年 $^{[28]}$，其直径约为太阳的 2.8 倍，光度约为太阳的$47$倍$^{[44]}$。小熊座$\varepsilon$星是一个三星系统$^{[45]}$，其合成平均视星等为4.22 $^{[46]}$。其主星为一颗光谱型为 G5III 的黄色巨星$^{[46]}$，同时也是一颗猎犬座 RS 型变星。该系统为一对光谱双星，伴星距离约 $0.36$ AU，另有第三颗恒星——一颗光谱型为 K0 的橙色主序星——位于约 $8100$ AU 处 $^{[45]}$。

靠近北极星的位置有小熊座$\lambda$星，这是一颗光谱型为 M1III 的红巨星。它是一颗半规则变星，视星等在$6.35$与$6.45$之间变化$^{[47]}$。由于该星座位于高北天区，其变星可全年观测：红巨星小熊座 R 星是一颗半规则变星，在$328$天内视星等由$8.5$变至$11.5$；而小熊座 S 星则是一颗长周期变星，在$331$天内视星等介于 8.0 与 11 之间$^{[48]}$。位于 Kochab 与 Pherkad 以南、朝向天龙座方向的是小熊座 RR 星$^{[33]}$，这是一颗光谱型为 M5III 的红巨星，同样属于半规则变星，其视星等在 43.3 天内于 4.44 至 4.85 之间变化 $^{[49]}$。小熊座 T 星是另一颗红巨星变星，其状态曾发生显著变化——从视星等在7.8至 15、周期310--315天的长周期（米拉型）变星，转变为半规则变星$^{[50]}$。该恒星被认为经历了一次“壳层氦闪”，即围绕恒星核心的氦壳层达到临界质量并发生点火，这一事件由其在$^{[1979]}$年变光性质的突然改变所标志$^{[51]}$。小熊座 Z 星是一颗微弱的变星，其在$^{[1992]}$年突然变暗$6$个星等，被确认属于一种极为罕见的恒星类别——北冕座 R 型变星$^{[52]}$。

食变星是由于一颗恒星从另一颗前方经过而导致亮度变化的恒星系统，而非源于自身光度变化。小熊座 W 星即为此类系统，其视星等在 1.7 天内于 8.51 与$9.59$ 之间变化$^{[53]}$。该系统的合成光谱型为 A2V，但两颗组成恒星的质量尚不明确。$^{[1973]}$年轨道周期的轻微变化表明该多星系统中可能存在第三个成员——很可能是一颗红矮星——其轨道周期为 $62.2\pm3.9$年$^{[54]}$。小熊座 RU 星是另一例，其视星等在 0.52 天内介于 10 与 10.66 之间 $^{[55]}$。这是一个半分离系统，因为次星充满了其洛希瓣，并正在向主星传递物质$^{[56]}$。

小熊座 RW 星是一个灾变变星系统，其在$^{[1956]}$年爆发为新星，亮度达到$6$等。到$^{[2003]}$年时，其亮度仍比基线值高出 2 个星等，并以每年 0.02 星等的速度逐渐变暗。其距离被计算为$5000\pm800$秒差距（约$16300$光年），这意味着它位于银河系晕之中$^{[57]}$。
“Calvera”这一绰号取自电影《豪勇七蛟龙》中的反派人物，用以指代 ROSAT 全天巡天亮源星表（RASS/BSC）中编号为 1RXS J141256.0+792204 的一处 X 射线源 $^{[58]}$。该天体已被确认是一颗孤立的中子星，是距离地球最近的此类天体之一 $^{[59]}$。小熊座中还包含两颗性质颇为神秘的白矮星。于 $^{[2011\text{-}01\text{-}27]}$ 记录的 H1504+65 是一颗极为暗弱（视星等 $15.9$）的恒星，却拥有目前已知白矮星中最高的表面温度——$200000\,\mathrm{K}$。其大气层由约一半的碳、一半的氧以及 $2\%$ 的氖组成，几乎不含氢和氦，这种成分目前尚无法用现有的恒星演化模型加以解释 $^{[60]}$。WD 1337+705 是一颗温度较低的白矮星，其光谱中含有镁和硅的特征线，这暗示其可能存在伴星或环绕恒星的碎屑盘，但至今尚未发现直接证据$^{[61]}$。WISE 1506+7027 是一颗光谱型为 T6 的棕矮星，距离地球仅约 $11.1^{+2.3}_{-1.3}$ 光年$^{[62]}$。该天体视星等为 14，于 $^{[2011]}$ 年由广域红外巡天探测器（WISE）发现$^{[63]}$。

除 Kochab 之外，还有另外三个恒星系统被发现拥有行星。11 小熊座是一颗光谱型为 K4III 的橙色巨星，其质量约为太阳的 1.8 倍。其年龄约为 15 亿年，自作为 A 型主序星以来已逐渐冷却并发生膨胀。该星距离地球约 390 光年，视星等为 5.04。$^{[2009]}$ 年，人们在其周围发现了一颗质量约为木星 11 倍的行星，其公转周期为 516 天 $^{[64]}$。HD 120084 是另一颗演化恒星，为光谱型 G7III 的黄色巨星，质量约为太阳的 2.4 倍。通过 $^{[2013]}$ 年对该恒星径向速度的精确测量，人们发现其拥有一颗质量约为木星 4.5 倍的行星，该行星具有极高的轨道偏心率（$e=0.66$），是已知偏心率最大的行星轨道之一 $^{[65]}$。HD 150706 是一颗光谱型为 G0V 的类太阳恒星，距离太阳系约 89 光年。最初曾认为其在 $0.6\,\mathrm{AU}$ 处拥有一颗质量与木星相当的行星，但这一说法在 $^{[2007]}$ 年被否定$^{[66]}$。随后于$^{[2012]}$ 年发表的进一步研究表明，该恒星确实存在一颗伴星，其质量约为木星的 2.7 倍，公转周期约为 $16$ 年，轨道半径约为 $6.8\,\mathrm{AU}$,$^{[67]}$。
\subsubsection{深空天体}
